% GENERATED FILE. Edit canonical lessons, then run npm run generate:books.
% canonical-source-hash: fnv1a64:45446987b1560e0f
% canonical-lessons: GU-C345-inkar, GU-C345-khulaso, GU-C345-cokhvat, GU-C345-ganavvu, GU-C345-palatvu

\chapter{Explaining Yourself}
\label{ch:gu-explaining-yourself}
\hlchaptermodality{345}

\begin{chapteropening}
\textbf{By the end of this chapter:} \emph{I can explain what I meant and set the record straight.}

Complete the last lesson of chapter 345: I can explain what I meant and set the record straight.
\end{chapteropening}

\section[\texorpdfstring{\gu{ઇનકાર} (inkār)}{ઇનકાર (inkār)}]{\gu{ઇનકાર} (inkār) — a refusal}

\label{lesson:GU-C345-inkar}



\begin{quote}
Before the new one: say the Gujarati for uncertainty, then the Gujarati for to indicate.
\end{quote}

\subsection*{You'll want to know: \gu{ઇનકાર}}
\textbf{\gu{ઇનકાર}} — \emph{inkār} — \textquotedblleft{}a refusal\textquotedblright{}.

\begin{quote}
\gu{દુકાનદારે} \gu{પૈસા} \gu{પાછા} \gu{આપવાનો} \gu{ઇનકાર} \gu{કર્યો}.
\end{quote}

\begin{quote}
\emph{dukāndāre paisā pāchā āpvāno inkār karyo.}
\end{quote}

\begin{quote}
The shopkeeper refused to give the money back.
\end{quote}

A masculine noun.

\begin{cousinweb}
\textbf{\gu{ઇનકાર}} (\emph{inkār}) is usually traced to Arabic \emph{inkār}, \textquotedblleft{}denial\textquotedblright{}, through Persian.
\end{cousinweb}

\begin{grammarlens}[title={What you've built}]
One more word for giving reasons and opinions.
\end{grammarlens}

\subsection*{Your turn}
\begin{itemize}
\raggedright
  \item \emph{Say it:} \emph{inkār}
  \item \emph{Say it:} \emph{inkār}, once more
  \item \emph{Say it:} \emph{dukāndāre paisā pāchā āpvāno inkār karyo.} — The shopkeeper refused to give the money back.
  \item \emph{Recall:} say the Gujarati for uncertainty, then the Gujarati for to indicate, then say \emph{inkār} again
\end{itemize}

\begin{tcolorbox}[breakable,colback=teal!4,colframe=teal!35!black,title={Before you move on}]
What does \gu{ઇનકાર} mean? (\textquotedblleft{}A refusal\textquotedblright{}.) Say it once more.
\end{tcolorbox}
\section[\texorpdfstring{\gu{ખુલાસો} (khulāso)}{ખુલાસો (khulāso)}]{\gu{ખુલાસો} (khulāso) — an explanation (of one's conduct)}

\label{lesson:GU-C345-khulaso}



\begin{quote}
Before the new one: say the Gujarati for to indicate, then the Gujarati for a refusal.
\end{quote}

\subsection*{You'll want to know: \gu{ખુલાસો}}
\textbf{\gu{ખુલાસો}} — \emph{khulāso} — \textquotedblleft{}an explanation (of one's conduct)\textquotedblright{}.

\begin{quote}
\gu{તેણે} \gu{મોડા} \gu{આવવાનો} \gu{ખુલાસો} \gu{કર્યો}.
\end{quote}

\begin{quote}
\emph{teṇe moḍā āvvāno khulāso karyo.}
\end{quote}

\begin{quote}
He explained why he was late.
\end{quote}

A masculine noun.

\begin{cousinweb}
\textbf{\gu{ખુલાસો}} (\emph{khulāso}) is usually traced to Arabic \emph{khulāṣa}, \textquotedblleft{}an extract, the gist\textquotedblright{}, through Persian; Gujarati uses it for clearing up a matter.
\end{cousinweb}

\begin{grammarlens}[title={What you've built}]
One more word for giving reasons and opinions.
\end{grammarlens}

\subsection*{Your turn}
\begin{itemize}
\raggedright
  \item \emph{Say it:} \emph{khulāso}
  \item \emph{Say it:} \emph{khulāso}, once more
  \item \emph{Say it:} \emph{teṇe moḍā āvvāno khulāso karyo.} — He explained why he was late.
  \item \emph{Recall:} say the Gujarati for to indicate, then the Gujarati for a refusal, then say \emph{khulāso} again
\end{itemize}

\begin{tcolorbox}[breakable,colback=teal!4,colframe=teal!35!black,title={Before you move on}]
What does \gu{ખુલાસો} mean? (\textquotedblleft{}An explanation (of one's conduct)\textquotedblright{}.) Say it once more.
\end{tcolorbox}
\section[\texorpdfstring{\gu{ચોખવટ} (cokhvaṭ)}{ચોખવટ (cokhvaṭ)}]{\gu{ચોખવટ} (cokhvaṭ) — a clear understanding (on a point)}

\label{lesson:GU-C345-cokhvat}



\begin{quote}
Before the new one: say the Gujarati for a refusal, then the Gujarati for an explanation.
\end{quote}

\subsection*{You'll want to know: \gu{ચોખવટ}}
\textbf{\gu{ચોખવટ}} — \emph{cokhvaṭ} — \textquotedblleft{}a clear understanding (on a point)\textquotedblright{}.

\begin{quote}
\gu{પૈસાની} \gu{ચોખવટ} \gu{પહેલેથી} \gu{કરવી} \gu{જોઈએ}.
\end{quote}

\begin{quote}
\emph{paisānī cokhvaṭ pahelethī karvī joīe.}
\end{quote}

\begin{quote}
Money matters should be made clear beforehand.
\end{quote}

A feminine noun.

\begin{cousinweb}
\textbf{\gu{ચોખવટ}} (\emph{cokhvaṭ}) is built on \emph{cokh}, \textquotedblleft{}clean, clear\textquotedblright{}, the base of \textbf{\gu{ચોખ્ખું}} (\emph{cokhkhũ}), which goes back to Sanskrit \emph{cokṣa}, \textquotedblleft{}pure, clean\textquotedblright{} (CDIAL 4918): a matter made clean and clear.
\end{cousinweb}

\begin{grammarlens}[title={What you've built}]
One more word for giving reasons and opinions.
\end{grammarlens}

\subsection*{Your turn}
\begin{itemize}
\raggedright
  \item \emph{Say it:} \emph{cokhvaṭ}
  \item \emph{Say it:} \emph{cokhvaṭ}, once more
  \item \emph{Say it:} \emph{paisānī cokhvaṭ pahelethī karvī joīe.} — Money matters should be made clear beforehand.
  \item \emph{Recall:} say the Gujarati for a refusal, then the Gujarati for an explanation, then say \emph{cokhvaṭ} again
\end{itemize}

\begin{tcolorbox}[breakable,colback=teal!4,colframe=teal!35!black,title={Before you move on}]
What does \gu{ચોખવટ} mean? (\textquotedblleft{}A clear understanding (on a point)\textquotedblright{}.) Say it once more.
\end{tcolorbox}
\section[\texorpdfstring{\gu{ગણાવવું} (gaṇāvvũ)}{ગણાવવું (gaṇāvvũ)}]{\gu{ગણાવવું} (gaṇāvvũ) — to list (reasons, one by one)}

\label{lesson:GU-C345-ganavvu}



\begin{quote}
Before the new one: say the Gujarati for an explanation, then the Gujarati for a clear understanding.
\end{quote}

\subsection*{You'll want to know: \gu{ગણાવવું}}
\textbf{\gu{ગણાવવું}} — \emph{gaṇāvvũ} — \textquotedblleft{}to list (reasons, one by one)\textquotedblright{}.

\begin{quote}
\gu{અરુણે} \gu{નોકરી} \gu{છોડવાનાં} \gu{ત્રણ} \gu{કારણો} \gu{ગણાવ્યાં}.
\end{quote}

\begin{quote}
\emph{Aruṇe nokrī choḍvānā̃ traṇ kāraṇo gaṇāvyā̃.}
\end{quote}

\begin{quote}
Arun gave three reasons for leaving his job.
\end{quote}

Past, with \gu{તેણે} (he or she …): \gu{તેણે} \gu{બધાં} \gu{કારણો} \gu{ગણાવ્યાં}, \emph{teṇe badhā̃ kāraṇo gaṇāvyā̃}.

\begin{cousinweb}
\textbf{\gu{ગણાવવું}} (\emph{gaṇāvvũ}) is the causative of \textbf{\gu{ગણવું}} (\emph{gaṇvũ}), \textquotedblleft{}to count\textquotedblright{}, from Sanskrit \emph{gaṇayati}, \textquotedblleft{}counts\textquotedblright{} (CDIAL 3993): to count things out for someone.
\end{cousinweb}

\begin{grammarlens}[title={What you've built}]
One more word for giving reasons and opinions.
\end{grammarlens}

\subsection*{Your turn}
\begin{itemize}
\raggedright
  \item \emph{Say it:} \emph{gaṇāvvũ}
  \item \emph{Say it:} \emph{gaṇāvvũ}, once more
  \item \emph{Say it:} \emph{Aruṇe nokrī choḍvānā̃ traṇ kāraṇo gaṇāvyā̃.} — Arun gave three reasons for leaving his job.
  \item \emph{Recall:} say the Gujarati for an explanation, then the Gujarati for a clear understanding, then say \emph{gaṇāvvũ} again
\end{itemize}

\begin{tcolorbox}[breakable,colback=teal!4,colframe=teal!35!black,title={Before you move on}]
What does \gu{ગણાવવું} mean? (\textquotedblleft{}To list (reasons, one by one)\textquotedblright{}.) Say it once more.
\end{tcolorbox}
\section[\texorpdfstring{\gu{પલટવું} (palaṭvũ)}{પલટવું (palaṭvũ)}]{\gu{પલટવું} (palaṭvũ) — to change one's stand}

\label{lesson:GU-C345-palatvu}



\begin{quote}
Before the new one: say the Gujarati for a clear understanding, then the Gujarati for to list.
\end{quote}

\subsection*{You'll want to know: \gu{પલટવું}}
\textbf{\gu{પલટવું}} — \emph{palaṭvũ} — \textquotedblleft{}to change one's stand\textquotedblright{}.

\begin{quote}
\gu{કાલે} \gu{તેણે} \gu{હા} \gu{કહી} \gu{અને} \gu{આજે} \gu{તે} \gu{પલટી} \gu{ગયો}.
\end{quote}

\begin{quote}
\emph{kāle teṇe hā kahī ane āje te palṭī gayo.}
\end{quote}

\begin{quote}
Yesterday he said yes, and today he has gone back on it.
\end{quote}

Past (he …): \gu{તે} \gu{બીજા} \gu{દિવસે} \gu{પલટી} \gu{ગયો}, \emph{te bījā divse palṭī gayo}.

\begin{cousinweb}
\textbf{\gu{પલટવું}} (\emph{palaṭvũ}) CDIAL lists the Gujarati verb, \textquotedblleft{}to change\textquotedblright{}, under a reconstructed \emph{pallaṭṭ} (CDIAL 7968).
\end{cousinweb}

\begin{grammarlens}[title={What you've built}]
One more word for giving reasons and opinions.
\end{grammarlens}

\subsection*{Your turn}
\begin{itemize}
\raggedright
  \item \emph{Say it:} \emph{palaṭvũ}
  \item \emph{Say it:} \emph{palaṭvũ}, once more
  \item \emph{Say it:} \emph{kāle teṇe hā kahī ane āje te palṭī gayo.} — Yesterday he said yes, and today he has gone back on it.
  \item \emph{Recall:} say the Gujarati for a clear understanding, then the Gujarati for to list, then say \emph{palaṭvũ} again
\end{itemize}

\begin{tcolorbox}[breakable,colback=teal!4,colframe=teal!35!black,title={Before you move on}]
What does \gu{પલટવું} mean? (\textquotedblleft{}To change one's stand\textquotedblright{}.) Say it once more.
\end{tcolorbox}
